\section{Model}\label{sec:model}

We consider two retail destinations located at the endpoints of a unit interval. Destination $L$ is located at $0$ and destination $R$ at $1$. A unit mass of shoppers is uniformly distributed on $z\in[0,1]$. The endpoint normalization follows the quadratic spatial-competition benchmark of \citet{daspremont1979}. Retailers have a common constant marginal cost $c$ and choose prices $p_L$ and $p_R$ simultaneously.

The distinguishing feature is that access quality is supplied by a third-party transport operator. The operator has a fixed total service resource $F>0$ that can be reallocated across the two directions. In addition to shoppers, there is an exogenous mass $M\geq0$ of background passengers traveling toward the $L$ direction. The background demand can represent commuting, education, or other travel that is independent of the retail price game.

\subsection{Shopper choice}

Let $t>0$ be the coefficient on quadratic spatial mismatch and let $w>0$ scale waiting disutility. Shoppers have inelastic unit demand over the two destinations: each shopper chooses either $L$ or $R$, and no outside option is modeled. The common gross value $v$ is therefore a welfare normalization and drops out of destination choice. If the operator supplies directional service levels $f_L$ and $f_R$, the utilities of a shopper located at $z$ are
\begin{align}
U_L(z)&=v-p_L-tz^2-\frac{w}{f_L},\\
U_R(z)&=v-p_R-t(1-z)^2-\frac{w}{f_R}.
\end{align}
Let $x\in[0,1]$ denote the shopper share served by $L$. The marginal shopper satisfies
\begin{equation}
 p_L-p_R+t(2x-1)+\frac{w}{f_L}-\frac{w}{f_R}=0.
 \label{eq:shopper_indifference_raw}
\end{equation}
Retail profits are
\begin{equation}
\pi_L=(p_L-c)x,
\qquad
\pi_R=(p_R-c)(1-x).
\label{eq:profits}
\end{equation}

\subsection{Third-party service allocation}

Conditional on the shopper share $x$, total directional passenger demands are
\begin{equation}
D_L=M+x,
\qquad
D_R=1-x.
\label{eq:directional_demand}
\end{equation}
The operator allocates the fixed resource $F$ to minimize aggregate waiting disutility subject to a minimum service requirement:
\begin{align}
\min_{f_L,f_R}\quad &w\left(\frac{M+x}{f_L}+\frac{1-x}{f_R}\right),\label{eq:operator_problem}\\
\text{s.t.}\quad & f_L+f_R=F,\qquad f_L,f_R\geq \underline f.
\end{align}
Define the minimum service share
\begin{equation}
q\equiv \frac{\underline f}{F}\in(0,1/2).
\end{equation}
Absent the floor, the unique cost-minimizing share of service allocated to direction $L$ is
\begin{equation}
s^u(x,M)
=\frac{\sqrt{M+x}}{\sqrt{M+x}+\sqrt{1-x}}.
\label{eq:unconstrained_share}
\end{equation}
With the service obligation, the implemented service share is the clipped rule
\begin{equation}
s_q(x,M)=\min\{1-q,\max\{q,s^u(x,M)\}\},
\label{eq:clipped_share}
\end{equation}
so that
\begin{equation}
f_L=Fs_q(x,M),
\qquad
f_R=F[1-s_q(x,M)].
\end{equation}
Let
\begin{equation}
A\equiv \frac{w}{F},
\qquad
h_q(x,M)\equiv\frac{1}{s_q(x,M)}-\frac{1}{1-s_q(x,M)}.
\end{equation}
Substituting the operator response into \eqref{eq:shopper_indifference_raw}, the fulfilled-expectations shopper condition becomes
\begin{equation}
G_q(x;p_L,p_R)
\equiv p_L-p_R+t(2x-1)+A h_q(x,M)=0.
\label{eq:Gq}
\end{equation}
When the service floor is slack,
\begin{equation}
h_q(x,M)=H(x,M)
\equiv \frac{1-M-2x}{\sqrt{(M+x)(1-x)}}.
\label{eq:H}
\end{equation}
The fixed-resource constraint is central: an increase in service toward one direction necessarily withdraws service from the other.

\subsection{Timing and equilibrium concept}

Retailers first choose prices simultaneously. Given a price pair, the shopper share and operator allocation are then determined jointly as a fulfilled-expectations continuation. Formally, for any candidate share $x$ the operator solves \eqref{eq:operator_problem}; shoppers choose between the two destinations taking the resulting service levels as given; and a continuation is a share that is consistent with those choices through \eqref{eq:Gq}, with the usual corner allocation if its zero lies outside $[0,1]$. This is a fixed-point description of the continuation, not a claim that the operator literally observes realized shopping demand before shoppers choose.

Retailers anticipate this continuation mapping when setting prices. We study pure-strategy Nash equilibrium in retail prices with the shopper/operator continuation solved globally. A pair of local price first-order conditions is therefore not sufficient for equilibrium.

The canonical parameter domain is
\begin{equation}
t>0,\quad F>0,\quad w>0,\quad M\geq0,\quad q\in(0,1/2),
\end{equation}
with inelastic unit shopper demand. The main result is an existence result on a nonempty open subset of this domain, not a claim that the same equilibrium structure obtains for all admissible primitives.
